\documentclass[a4paper,12pt]{article}
\usepackage[utf8]{inputenc}
\usepackage[T1]{fontenc}
\usepackage[margin=2.5cm]{geometry}
\usepackage{ctex} % 支持中文
\usepackage{booktabs} % 专业表格线
\usepackage{array} % 表格列格式
\usepackage{graphicx} % 图片支持
\usepackage{titlesec} % 标题格式
\usepackage{xcolor} % 颜色支持
\usepackage{enumitem} % 列表优化
\usepackage{float} % 浮动体控制
\usepackage{setspace} % 行距控制

% 设置标题格式
\titleformat{\section}{\Large\bfseries\color{blue!40!black}}{}{0.5em}{}
\titleformat{\subsection}{\large\bfseries\color{gray!80!black}}{}{0.5em}{}

% 列表设置
\setlist[itemize]{leftmargin=*}

% 行距设置
\onehalfspacing

% 元数据
\title{\textbf{2024年度环境、社会及管治（ESG）报告节选}}
\author{公司董事会办公室}
\date{2024年3月}

\begin{document}

\maketitle

% 依据指令，已移除摘要和目录

\section{1. 前言：战略协同与可持续发展框架综述}

2024年，被界定为海南自由贸易港封关运作准备工作的攻坚之年，亦是本公司（以下简称“公司”）作为区域性基础设施建设与临空产业运营主体，深化高质量发展路径的关键时期。鉴于公司在“一主一次，两支一货运”区域机场群协同发展战略中的核心地位，以及其对于国家对外开放战略布局的实质性支撑作用，本报告旨在系统性地披露公司在环境、社会及管治（ESG）维度的管理效能与绩效表现。公司致力于通过建立健全的内部控制体系与可持续发展治理架构，确保“外防输入、内保畅通”运营目标的达成，并在此基础上，通过对环境足迹的精细化管控及社会价值的多元化创造，为包括股东在内的所有利益相关方构建长期稳健的价值回报机制。

\section{2. 环境范畴（Environmental）：碳排放管控与资源集约化利用}

在应对全球气候变化及响应国家“双碳”战略目标的宏观背景下，公司已将环境管理纳入企业核心战略范畴。报告期内，公司依据适用的法律法规及行业标准，实施了严格的温室气体排放核算与能源审计程序，旨在通过技术迭代与管理优化，实现运营碳强度的持续下降。

\subsection{2.1 温室气体排放与资源消耗量化绩效}

依据ISO 14064标准及相关行业温室气体核算指南，公司对2024年度的碳排放数据进行了系统性盘查。相关核心环境绩效指标如下表所示：

\begin{table}[H]
    \centering
    \caption{2024年度环境绩效关键指标统计表}
    \label{tab:env_kpi}
    \renewcommand{\arraystretch}{1.3}
    \begin{tabular}{p{9cm} p{3cm} p{3cm}}
        \toprule
        \textbf{环境绩效指标} & \textbf{统计数值} & \textbf{计量单位} \\
        \midrule
        \multicolumn{3}{l}{\textit{温室气体排放（Greenhouse Gas Emissions）}} \\
        温室气体排放总量 & 51,518.71 & tCO$_2$e \\
        \quad - 范围1（直接排放：固定源与移动源燃烧） & 11,214.39 & tCO$_2$e \\
        \quad - 范围2（间接排放：外购电力消费） & 40,304.32 & tCO$_2$e \\
        单位营收碳排放强度 & 0.12 & tCO$_2$e/万元 \\
        \midrule
        \multicolumn{3}{l}{\textit{能源与资源消耗（Energy \& Resource Consumption）}} \\
        综合能源消耗量 & 13,538.22 & 吨标准煤 \\
        外购电力消耗量 & 7,511.05 & 万千瓦时 \\
        移动源汽油消耗量 & 19.10 & 万升 \\
        \midrule
        \multicolumn{3}{l}{\textit{水资源管理与废弃物处置（Water \& Waste Management）}} \\
        水资源消耗总量 & 145.94 & 万吨 \\
        中水回用量（再生水利用） & 13.61 & 万吨 \\
        无害废弃物产生量 & 6,230.58 & 吨 \\
        有害废弃物产生量 & 0.58 & 吨 \\
        废弃物合规处置率 & 100 & \% \\
        \bottomrule
    \end{tabular}
\end{table}

\subsection{2.2 绿色机场建设路径与减排技术应用}

为达成“绿色机场”的建设目标，公司采取了一系列旨在降低航空器地面运行能耗及提升可再生能源利用率的技术干预措施。

\begin{itemize}
    \item \textbf{辅助动力装置（APU）替代设施的部署与效能}：
    \begin{itemize}
        \item \textbf{实施机制}：通过在停机位部署地面电源装置（GPU）及地面空调组件（PCA），以替代航空器在地面停泊期间对机载辅助动力装置（APU）的依赖，从而显著削减航空燃油消耗及相应的温室气体与噪声排放。
        \item \textbf{设施规模}：截至报告期末，旗下核心门户机场已投入运营的APU替代设施共计31套，涵盖18个近机位与13个远机位。
        \item \textbf{减排绩效}：经测算，该项目的实施在2024年度实现了约 \textbf{25,382} 吨二氧化碳当量的减排效应。
        \item \textbf{认证获取}：基于上述成效，旗下核心门户机场被授予“双碳机场”二星级认证。
    \end{itemize}
    
    \item \textbf{分布式光伏发电系统的整合应用}：
    \begin{itemize}
        \item \textbf{商业建筑应用}：旗下核心商业综合体实施了屋顶分布式光伏项目，其年均发电量达 450 万千瓦时，占该建筑体总电力需求的约10\%，折算年度碳减排量为 2,263.59 吨。
        \item \textbf{产业园区应用}：旗下临空产业园利用闲置屋顶资源，完成了装机容量为 8.6MWp 的光伏发电系统建设，预计年均碳减排贡献可达 25.75 万吨。
    \end{itemize}

    \item \textbf{场内车辆电动化转型}：
    \begin{itemize}
        \item 截至2024年末，核心机场区域内的新能源车辆保有量已提升至 221 辆，占车辆总数的 43.6\%；同期配套建设充电桩设施约 100 个。
        \item 该项车辆能源结构调整措施在报告期内实现了 761 吨的碳排放削减。
    \end{itemize}
\end{itemize}

\subsection{2.3 绿色建筑认证体系}

在建筑资产的开发与运营全生命周期中，公司严格遵循国际及国内绿色建筑评价标准：
\begin{itemize}
    \item 省会城市在建地标性项目“海南中心”（双子塔）已成功获取 \textbf{LEED CS 金级预认证}，体现了其在建筑节能与环境设计领域的卓越表现。
    \item 位于某新区的F07项目经评估符合并获得了绿色建筑一星级标准认证。
\end{itemize}

\section{3. 社会范畴（Social）：运营韧性构建与利益相关方价值共创}

公司在确保航空安全运营底线的基础上，致力于通过服务流程的再造与应急响应机制的强化，提升运营韧性；同时，坚持多元包容的人力资源政策，并积极履行企业公民责任。

\subsection{3.1 核心运营指标与业务恢复情况}

报告期内，受宏观经济复苏及旅游市场需求反弹的驱动，公司主要运营指标呈现稳健增长态势。

\begin{table}[H]
    \centering
    \caption{2024年度核心运营数据概览}
    \label{tab:ops_data}
    \renewcommand{\arraystretch}{1.3}
    \begin{tabular}{p{8cm} p{5cm}}
        \toprule
        \textbf{运营指标项目} & \textbf{统计数据} \\
        \midrule
        旗下控股及管理输出机场总数 & 9 家 \\
        年度总体起降架次 & 16.11 万架次 \\
        年度旅客吞吐量 & 2,462.62 万人次 \\
        核心门户机场行业位次 & 重回“2000万级”梯队（全国第29位） \\
        重点区域机场网络拓展 & 开通至吉隆坡的首条国际客运航线 \\
        \bottomrule
    \end{tabular}
\end{table}

\subsection{3.2 安全管理体系（SMS）与极端天气应急响应}

公司建立并维持了一套严谨的安全管理体系，并针对极端气象灾害制定了专项应急预案，以确保业务连续性。

\begin{itemize}
    \item \textbf{超强台风“摩羯”的应急响应与处置}：
    \begin{itemize}
        \item \textbf{预警与防御}：启动“防汛防风专项应急预案”，运管委实施了包括航班动态调整及航空器系留加固在内的预防性措施。
        \item \textbf{灾后恢复}：台风过境后，迅速组织设施修复与适航检查。9月3日至7日期间，核心机场在确保安全的前提下迅速恢复运行，累计保障航班 \textbf{1,300} 架次，有效保障了区域交通枢纽的畅通。
    \end{itemize}
    
    \item \textbf{数字化安全管控平台的应用}：
    \begin{itemize}
        \item 上线“安全智控”平台，实现安全管理要素的数字化与流程化。
        \item 平台关联安全制度手册条款 9,626 项，全年流转并处置电子检查单 31 万份。
    \end{itemize}

    \item \textbf{年度安全绩效}：
    \begin{itemize}
        \item 一般征候万架次率：0
        \item 责任事故发生率：0
        \item 安全培训全员覆盖率：100\% \\
    \end{itemize}
\end{itemize}

\subsection{3.3 服务流程优化与创新}

\begin{itemize}
    \item \textbf{查验流程集约化}：实施“一机双屏”海关与安检融合查验模式，通过图像共享技术实现旅客“一次过检”，显著提升通关效能。
    \item \textbf{入境服务便利化}：在核心机场设立“一站式综合服务中心”，集成支付绑定、通信服务及文旅咨询功能，解决外籍旅客入境服务的“最后一公里”问题。
\end{itemize}

\subsection{3.4 人力资本与社会投入}

\begin{table}[H]
    \centering
    \caption{2024年度人力资源与社会贡献统计}
    \label{tab:social_stats}
    \renewcommand{\arraystretch}{1.3}
    \begin{tabular}{p{7cm} p{6cm}}
        \toprule
        \textbf{指标维度} & \textbf{统计数据} \\
        \midrule
        \textbf{人力资源构成} & \\
        期末员工总数 & 10,611 人 \\
        女性员工占比 & 38.9\% (4,133人) \\
        少数民族员工人数 & 179 人 \\
        \midrule
        \textbf{培训与发展} & \\
        年度培训投入总额 & 1,119.54 万元 \\
        人均培训时长 & 105.88 小时 \\
        \midrule
        \textbf{社会公益} & \\
        乡村振兴资金投入 & 190 万元 \\
        志愿服务参与人次 & >10,000 人次 \\
        志愿服务总时长 & 500,000 小时 \\
        对外捐赠资金 & 15.48 万元 \\
        \bottomrule
    \end{tabular}
\end{table}

\section{4. 管治范畴（Governance）：治理架构与合规监督体系}

公司致力于构建符合上市公司监管要求的高标准治理架构，通过强化内部控制与合规管理，保障决策的科学性与透明度。

\subsection{4.1 治理结构与信息披露}

\begin{itemize}
    \item \textbf{董事会构成}：董事会由9名成员组成，其中包含3名独立董事，符合独立性监管要求；女性董事占比为 22\%。
    \item \textbf{会议与决策}：报告期内，共召开股东大会5次，董事会会议13次，确保了重大事项决策程序的合法合规。
    \item \textbf{信息披露合规}：全年发布各类公告70份，严格履行了信息披露义务。
\end{itemize}

\subsection{4.2 风险管控与商业道德}

公司实施了“三全三化”大监督体系，旨在实现对关键业务流程与风险点的全覆盖监控。

\begin{table}[H]
    \centering
    \caption{2024年度合规与风控管理绩效}
    \label{tab:gov_stats}
    \renewcommand{\arraystretch}{1.3}
    \begin{tabular}{p{9cm} p{4cm}}
        \toprule
        \textbf{合规管理指标} & \textbf{执行情况} \\
        \midrule
        风险点识别与管控数量 & 358 个 \\
        反腐败培训覆盖率（管理层） & 95\% \\
        反腐败培训覆盖率（员工层） & 90\% \\
        反腐败培训总时长 & 400 小时 \\
        供应商廉洁承诺书签署率 & 100\% \\
        工程类供应商遴选数量 & 23 家 \\
        \bottomrule
    \end{tabular}
\end{table}

\subsection{4.3 间接经济影响}

作为自贸港建设的积极参与者，公司通过参股投资方式布局离岛免税业务，对区域经济产生了实质性的间接经济影响。2024年度，公司间接参与的5家离岛免税店实现线下销售额约为 \textbf{52 亿元}，占全岛市场份额约 \textbf{17\%}，有效促进了消费回流与地方税收增长。

\end{document}